\documentclass[11pt,a4paper]{article}

\usepackage[utf8]{inputenc}
\usepackage[T1]{fontenc}
\usepackage[ngerman]{babel}

\usepackage{amsmath,amssymb,amsthm}
\usepackage{mathtools}
\usepackage[a4paper,margin=2.5cm]{geometry}
\usepackage{lmodern}
\usepackage{microtype}
\usepackage{enumitem}
\usepackage{hyperref}

\title{Geschichtete Radikalzerlegung, Sieben-Typen-Klassifikation\\
	und Radikalzerlegungsrest}
\author{Thomas Hoffbauer}
\date{\today}

% Theorem-Umgebungen
\newtheorem{satz}{Satz}
\newtheorem{lemma}[satz]{Lemma}
\newtheorem{korollar}[satz]{Korollar}
\theoremstyle{definition}
\newtheorem{definition}[satz]{Definition}
\newtheorem{beispiel}[satz]{Beispiel}
\theoremstyle{remark}
\newtheorem{bemerkung}[satz]{Bemerkung}

% Makros
\newcommand{\N}{\mathbb{N}}
\newcommand{\rad}{\operatorname{rad}}
\newcommand{\Typ}{\operatorname{Typ}}
\newcommand{\Rest}{\mathcal{R}}
\newcommand{\vp}{v_p}

\begin{document}
	\maketitle
	
	\begin{abstract}
		Wir formulieren eine geschichtete Radikalzerlegung natürlicher Zahlen nach Abspaltung der
		\(2\)- und \(3\)-adischen Anteile und führen darauf aufbauend eine reduzierte
		Sieben-Typen-Klassifikation ein. Diese Klassifikation basiert auf den vier Restklassen
		\(1,5,7,11 \pmod{12}\), interpretiert als Familien \(E,A,B,C\).
		Jede Radikalschicht besitzt genau einen reduzierten Typ in
		\[
		\{E,A,B,C,AB,AC,BC\},
		\]
		und zu jeder natürlichen Zahl existiert eine eindeutig bestimmte Folge solcher Typen,
		der sogenannte \emph{Radikalzerlegungsrest}.
	\end{abstract}
	
	\section{Einleitung}
	
	Der Fundamentalsatz der Arithmetik liefert für jede natürliche Zahl \(n\ge 1\) eine eindeutige
	Primfaktorzerlegung
	\[
	n=\prod_p p^{\vp(n)}.
	\]
	Diese Darstellung organisiert die Information primzahlweise, also vertikal entlang der einzelnen
	Primzahlen. Eine zweite, dazu äquivalente, aber strukturell anders geordnete Beschreibung ergibt
	sich, wenn man dieselbe Exponenteninformation horizontal nach \emph{Schichten} organisiert.
	
	Nach Abspaltung der Faktoren \(2\) und \(3\) betrachtet man nur noch den zu \(6\) teilerfremden Teil.
	Dessen Primteiler liegen modulo \(12\) in den vier Restklassen
	\[
	1,\;5,\;7,\;11.
	\]
	Diese werden im Folgenden als Familien \(E,A,B,C\) aufgefasst. Auf jeder Radikalschicht kann man
	dann die Primteiler paritätisch nach diesen Familien reduzieren. Daraus ergibt sich eine endliche
	Typenklassifikation mit genau sieben reduzierten Typen sowie eine eindeutig definierte Restfolge.
	
	\section{Geschichtete Radikalzerlegung}
	
	\begin{definition}[Radikal]
		Für \(m\in\N\) definiert man das Radikal durch
		\[
		\rad(m):=\prod_{p\mid m} p,
		\]
		wobei das Produkt über alle verschiedenen Primteiler von \(m\) läuft.
		Insbesondere ist \(\rad(m)\) quadratfrei.
	\end{definition}
	
	\begin{definition}[Abspaltung der Faktoren \(2\) und \(3\)]
		Für \(n\in\N\) setzen wir
		\[
		\alpha:=v_2(n),\qquad \beta:=v_3(n),\qquad m:=\frac{n}{2^\alpha 3^\beta}.
		\]
		Dann gilt \(\gcd(m,6)=1\).
	\end{definition}
	
	\begin{definition}[Radikalschichten]
		Sei
		\[
		m=\prod_p p^{v_p(m)}
		\]
		mit \(\gcd(m,6)=1\). Für
		\[
		k=1,\dots,h,\qquad h:=\max_{p\mid m} v_p(m),
		\]
		definieren wir die \(k\)-te Radikalschicht durch
		\[
		R_k(m):=\prod_{\substack{p\mid m\\ v_p(m)\ge k}} p.
		\]
		Jede Zahl \(R_k(m)\) ist quadratfrei und zu \(6\) teilerfremd.
	\end{definition}
	
	\begin{lemma}
		Für jedes \(m\in\N\) mit \(\gcd(m,6)=1\) gilt
		\[
		m=\prod_{k=1}^{h} R_k(m),
		\qquad
		h=\max_{p\mid m} v_p(m).
		\]
	\end{lemma}
	
	\begin{proof}
		Fixiere eine Primzahl \(p\mid m\). Dann tritt \(p\) genau in denjenigen Schichten \(R_k(m)\) auf,
		für die \(v_p(m)\ge k\) gilt. Das sind genau \(v_p(m)\) viele Schichten. Folglich ist der Exponent
		von \(p\) im Produkt \(\prod_{k=1}^{h} R_k(m)\) gleich \(v_p(m)\). Da dies für alle Primteiler gilt,
		haben beide Seiten dieselbe Primfaktorzerlegung.
	\end{proof}
	
	\begin{lemma}
		Die Folge der Radikalschichten ist divisibilitätsmonoton:
		\[
		R_{k+1}(m)\mid R_k(m)
		\qquad \text{für alle } k=1,\dots,h-1.
		\]
	\end{lemma}
	
	\begin{proof}
		Wenn \(p\mid R_{k+1}(m)\), dann ist \(v_p(m)\ge k+1\), also insbesondere \(v_p(m)\ge k\). Damit
		folgt \(p\mid R_k(m)\). Also teilt \(R_{k+1}(m)\) die Schicht \(R_k(m)\).
	\end{proof}
	
	\begin{satz}[Eindeutige geschichtete Radikalzerlegung]
		Zu jeder natürlichen Zahl \(n\in\N\) existieren eindeutig bestimmte Zahlen
		\[
		\alpha=v_2(n),\qquad \beta=v_3(n),
		\]
		sowie eindeutig bestimmte quadratfreie Zahlen
		\[
		R_1,\dots,R_h
		\]
		mit
		\[
		n=2^\alpha 3^\beta \prod_{k=1}^{h} R_k,
		\qquad
		R_{k+1}\mid R_k,
		\qquad
		\gcd(R_k,6)=1.
		\]
	\end{satz}
	
	\begin{proof}
		Die Eindeutigkeit von \(\alpha,\beta\) folgt aus der Primfaktorzerlegung. Der zu \(6\) teilerfremde
		Teil \(m=n/(2^\alpha 3^\beta)\) ist dadurch eindeutig festgelegt. Die Zahlen \(R_k=R_k(m)\) sind
		durch die Exponenten \(v_p(m)\) eindeutig bestimmt. Die beiden vorigen Lemmata liefern die Existenz
		und die gewünschte Struktur.
	\end{proof}
	
	\section{Die Familien \(E,A,B,C\)}
	
	\begin{definition}[Familien modulo \(12\)]
		Für Primzahlen \(p\ge 5\) definieren wir vier Familien:
		\[
		E:=\{p\text{ prim}: p\equiv 1 \pmod{12}\},
		\]
		\[
		A:=\{p\text{ prim}: p\equiv 5 \pmod{12}\},
		\]
		\[
		B:=\{p\text{ prim}: p\equiv 7 \pmod{12}\},
		\]
		\[
		C:=\{p\text{ prim}: p\equiv 11 \pmod{12}\}.
		\]
	\end{definition}
	
	Da jede Primzahl \(p\ge 5\) zu \(1,5,7\) oder \(11 \pmod{12}\) kongruent ist, liegt jeder Primteiler
	einer Zahl \(m\) mit \(\gcd(m,6)=1\) in genau einer dieser vier Familien.
	
	\begin{lemma}
		In der Einheitengruppe \((\mathbb{Z}/12\mathbb{Z})^\times\) gilt
		\[
		A^2=B^2=C^2=E,
		\]
		sowie
		\[
		AB=C,\qquad AC=B,\qquad BC=A,\qquad ABC=E.
		\]
	\end{lemma}
	
	\begin{proof}
		Dies ist eine direkte Rechnung mit den Restklassen
		\[
		5^2\equiv 1,\qquad 7^2\equiv 1,\qquad 11^2\equiv 1 \pmod{12},
		\]
		sowie
		\[
		5\cdot 7\equiv 11,\qquad 5\cdot 11\equiv 7,\qquad 7\cdot 11\equiv 5 \pmod{12}.
		\]
		Außerdem ist
		\[
		5\cdot 7\cdot 11\equiv 1 \pmod{12}.
		\]
	\end{proof}
	
	\section{Sieben reduzierte Typen}
	
	\begin{definition}[Paritätssignatur]
		Sei \(R\) eine quadratfreie Zahl mit \(\gcd(R,6)=1\). Wir definieren
		\[
		\varepsilon_A(R):=\#\{p\mid R:\ p\in A\}\bmod 2,
		\]
		\[
		\varepsilon_B(R):=\#\{p\mid R:\ p\in B\}\bmod 2,
		\]
		\[
		\varepsilon_C(R):=\#\{p\mid R:\ p\in C\}\bmod 2.
		\]
		Das Tripel
		\[
		(\varepsilon_A(R),\varepsilon_B(R),\varepsilon_C(R))\in\{0,1\}^3
		\]
		heißt Paritätssignatur von \(R\).
	\end{definition}
	
	Formal ergeben sich daraus acht Muster:
	\[
	\varnothing,\ A,\ B,\ C,\ AB,\ AC,\ BC,\ ABC.
	\]
	Wegen der Relation \(ABC=E\) werden jedoch \(\varnothing\) und \(ABC\) zum selben neutralen Typ
	zusammengefasst.
	
	\begin{definition}[Reduzierter Typ]
		Der \emph{reduzierte Typ} einer quadratfreien Zahl \(R\) mit \(\gcd(R,6)=1\) ist das eindeutig
		bestimmte Element
		\[
		\Typ(R)\in\{E,A,B,C,AB,AC,BC\},
		\]
		das durch die Paritätssignatur von \(R\) nach den Regeln
		\[
		(0,0,0)\mapsto E,\quad (1,0,0)\mapsto A,\quad (0,1,0)\mapsto B,\quad (0,0,1)\mapsto C,
		\]
		\[
		(1,1,0)\mapsto AB,\quad (1,0,1)\mapsto AC,\quad (0,1,1)\mapsto BC,\quad (1,1,1)\mapsto E
		\]
		festgelegt ist.
	\end{definition}
	
	\begin{satz}[Siebentypen-Klassifikation]
		Jede Radikalschicht \(R_k\) einer natürlichen Zahl \(n\) besitzt genau einen reduzierten Typ
		\[
		\Typ(R_k)\in\{E,A,B,C,AB,AC,BC\}.
		\]
		Insbesondere gibt es genau sieben reduzierte Typen von Radikalschichten.
	\end{satz}
	
	\begin{proof}
		Jede Radikalschicht \(R_k\) ist quadratfrei und zu \(6\) teilerfremd. Daher sind die Zahlen
		\(\varepsilon_A(R_k)\), \(\varepsilon_B(R_k)\), \(\varepsilon_C(R_k)\) eindeutig bestimmt. Dies
		liefert genau eines der acht formalen Muster. Da die beiden neutralen Muster \(\varnothing\) und
		\(ABC\) durch die Relation \(ABC=E\) denselben Typ \(E\) repräsentieren, bleiben genau sieben
		verschiedene reduzierte Typen übrig.
	\end{proof}
	
	\begin{bemerkung}
		Der Typ \(E\) ist kein bloßer Sonderfall, sondern ein vollwertiger Typ. Er umfasst insbesondere
		\begin{itemize}[itemsep=0.2em]
			\item einzelne Primzahlen \(p\equiv 1\pmod{12}\),
			\item Produkte solcher Primzahlen,
			\item sowie allgemein jede quadratfreie Zahl, deren Paritätssignatur neutral ist.
		\end{itemize}
	\end{bemerkung}
	
	\section{Der Radikalzerlegungsrest}
	
	\begin{definition}[Radikalzerlegungsrest]
		Sei
		\[
		n=2^\alpha 3^\beta \prod_{k=1}^{h} R_k
		\]
		die eindeutige geschichtete Radikalzerlegung von \(n\). Dann heißt die Folge
		\[
		\Rest(n):=\bigl(\Typ(R_1),\Typ(R_2),\dots,\Typ(R_h)\bigr)
		\]
		der \emph{Radikalzerlegungsrest} von \(n\).
	\end{definition}
	
	\begin{satz}[Existenz und Eindeutigkeit des Radikalzerlegungsrestes]
		Zu jeder natürlichen Zahl \(n\in\N\) existiert ein eindeutig bestimmter Radikalzerlegungsrest
		\[
		\Rest(n)\in \{E,A,B,C,AB,AC,BC\}^{<\infty}.
		\]
	\end{satz}
	
	\begin{proof}
		Die geschichtete Radikalzerlegung
		\[
		n=2^\alpha 3^\beta \prod_{k=1}^{h} R_k
		\]
		ist eindeutig. Für jede Schicht \(R_k\) ist der reduzierte Typ \(\Typ(R_k)\) eindeutig bestimmt.
		Damit ist auch die gesamte Folge \(\Rest(n)\) eindeutig.
	\end{proof}
	
	\begin{korollar}
		Jede natürliche Zahl \(n\) besitzt
		\begin{enumerate}[label=(\roman*),itemsep=0.2em]
			\item einen eindeutig bestimmten \(2\)-adischen Anteil \(2^{v_2(n)}\),
			\item einen eindeutig bestimmten \(3\)-adischen Anteil \(3^{v_3(n)}\),
			\item eine eindeutig bestimmte geschichtete Radikalzerlegung des zu \(6\) teilerfremden Restes,
			\item und einen eindeutig bestimmten Radikalzerlegungsrest \(\Rest(n)\).
		\end{enumerate}
	\end{korollar}
	
	\section{Beispiele}
	
	\begin{beispiel}
		Betrachte
		\[
		n=2^2\cdot 3\cdot 5^3\cdot 7^2\cdot 13.
		\]
		Dann ist
		\[
		m=\frac{n}{2^2\cdot 3}=5^3\cdot 7^2\cdot 13.
		\]
		Die Radikalschichten sind
		\[
		R_1=5\cdot 7\cdot 13,\qquad R_2=5\cdot 7,\qquad R_3=5.
		\]
		Da \(13\in E\), \(5\in A\) und \(7\in B\), folgt
		\[
		\Typ(R_1)=AB,\qquad \Typ(R_2)=AB,\qquad \Typ(R_3)=A.
		\]
		Also ist
		\[
		\Rest(n)=(AB,AB,A).
		\]
	\end{beispiel}
	
	\begin{beispiel}
		Für
		\[
		n=5\cdot 7\cdot 11
		\]
		gibt es nur eine Schicht
		\[
		R_1=5\cdot 7\cdot 11.
		\]
		Die Paritätssignatur ist formal \(ABC\), dies reduziert sich jedoch zu \(E\). Also
		\[
		\Rest(n)=(E).
		\]
	\end{beispiel}
	
	\begin{beispiel}
		Für die Primzahl
		\[
		n=13
		\]
		gilt \(13\equiv 1\pmod{12}\), also \(13\in E\). Damit ist
		\[
		\Rest(13)=(E).
		\]
		Für die Primzahl
		\[
		n=19
		\]
		gilt \(19\equiv 7\pmod{12}\), also \(19\in B\). Daher ist
		\[
		\Rest(19)=(B).
		\]
	\end{beispiel}
	
	\section{Schlussbemerkung}
	
	Die hier eingeführte Struktur trennt drei Ebenen voneinander:
	\begin{enumerate}[label=\arabic*.,itemsep=0.25em]
		\item die gewöhnliche Primfaktorzerlegung,
		\item die geschichtete Radikalzerlegung nach Exponentenstufen,
		\item die reduzierte Sieben-Typen-Signatur jeder Schicht.
	\end{enumerate}
	
	Damit erhält jede natürliche Zahl nicht nur eine eindeutige Primzerlegung, sondern zusätzlich eine
	eindeutige \emph{Signaturfolge} über dem Alphabet
	\[
	\{E,A,B,C,AB,AC,BC\}.
	\]
	Diese Folge ist der Radikalzerlegungsrest der Zahl und kann als reduzierte, schichtweise
	arithmetische Signatur interpretiert werden.
	
\end{document}